\scanpage{46}
On définit une catégorie $\unsure{R|\mathscr R}(\mathfrak X/\mathfrak X_0)$ ainsi.

$R(\mathfrak X)=$ catégorie des $\mathfrak X$-préschémas formels finis sur $\mathfrak X$.

\deleted{On définit la \unsure{catégor.|catégorie} des flèches $S$ [isom mod $\mathfrak X_0$] par $u:\mathfrak X'\mathbin{-}\mathfrak X''$ est un isom mod $\mathfrak X_0$ sss l'hom d'Algèbres $\mathcal A''\mathbin{-}\mathcal A'$ correspondant est un isom de Modules mod $X_0$. On vérifie qu'on obtient un calcul des fractions à droite [stable par composition, stable par changmt de base, égalisation]. La catégorie quotient est notée $R(\mathfrak X/\mathfrak X_0)$.}

\deleted{\uline{Exemple} Reprenons l'exemple précédent. Alors $R(\mathfrak X)\xleftarrow{\simeq}\unsure{\mathbb R|R}(X)$ par le théorème de comparaison, donc $R(\mathfrak X/\mathfrak X_0)$ est équivalente : $R(X/X_0)$, qu'on vérifie être équivalente à $R(X-X_0)$.}

\deleted{[$R(X/X_0)\mathbin{-}R(X-X_0)$, c'est \uline{\unsure{flt.|plt.|fth.} fidèle} comme on vérifie directement, c'est ess. surjectif comme on voit grâce à la version globale du Main Theorem p.ex. .....] .}

% slanted
\margin{\uline{à vérifier}}
% archive sheet number: 45
